\section{Scaling limit of the Mittag-Leffler ensemble}

In this section, we construct a fractional differential equation satisfied by the local kernel of the Mittag-Leffler ensemble with potentials $Q(\zeta)=|\zeta|^{2\lambda}-(2c/N)\log|\zeta|$, $\lambda>0$ and $c>-1$ (Proposition~\ref{Thm_bulkS}). As a consequence, we prove Theorem~\ref{Thm_ML integer}.

\subsection{Christoffel--Darboux type formula}
For the Mittag-Leffler potentials \eqref{Q ML}, the orthogonal norm $h_k$ is given by
\begin{gather*}
h_k=\int_\C |\zeta|^{2k} {\rm e}^{-N Q(\zeta)} \, {\rm d}A(\zeta)=\frac{\Gamma\big( \frac{k+c+1}{\lambda} \big)}{\lambda N^{\frac{k+c+1}{\lambda}} } .
\end{gather*}
Then by \eqref{skew op_rad}, we have
\begin{gather} \label{SOP ML}
	q_{2k+1}(\zeta)=\zeta^{2k+1}, \qquad q_{2k}(\zeta)=\zeta^{2k}+ \sum_{l=0}^{k-1} \frac{\zeta^{2l}}{N^{\frac{k-l}{\lambda}}} \prod_{j=0}^{k-l-1} \frac{ \Gamma\big( \frac{2l+2j+3+c}{\lambda} \big) }{ \Gamma\big( \frac{2l+2j+2+c}{\lambda} \big) } ,
\end{gather}
and
\begin{gather} \label{sk ML}
	s_k=2h_{2k+1}=\frac{2}{\lambda } \frac{\Gamma( \frac{2k+c+2}{\lambda} )}{N^{\frac{2k+c+2}{\lambda}} } .
\end{gather}
To describe the Christoffel--Darboux type formula, let us recall that the Caputo fractional derivative $D^\nu_z$ is given by
\begin{gather*}
	D^\nu_z f(z):=\frac{1}{\Gamma(n-\nu)} \int_{0}^{z} (z-u)^{n-\nu-1} f^{(n)}(u) \, {\rm d}u, \qquad n-1 < \nu <n, \quad n \in \mathbb{N},
\end{gather*}
see \cite[Section 2.4]{MR2218073}. Here $f^{(n)}$ is the usual $n$:th derivative.

Recall that $p=0$
%%
in the rescaling \eqref{zeta eta} here,
and $K_{\lambda,c}$ is given by \eqref{K ML}.

\begin{Proposition} \label{Thm_bulkS}
	For each $\lambda>0$ and $c>-1$, there exists a sequence of cocycles $(c_N)_{N=1}^\infty$ such that
	\begin{gather*}
			\lim_{N \to \infty} c_N\cdot K_N(z,w)=K_{\lambda,c}(z,w)
			={\rm e}^{ -\frac{ |z|^{2\lambda}+|w|^{2\lambda} }{\lambda} } \begin{pmatrix}
				\kappa_{\lambda,c}(z,w) & \kappa_{\lambda,c}(z,\bar{w}) \\
				\kappa_{\lambda,c}(\bar{z},w) & \kappa_{\lambda,c}(\bar{z},\bar{w})
			\end{pmatrix}
		\end{gather*}
	uniformly for $z$, $w$ in compact subsets of~$\C$, where
	\begin{gather*}
		\kappa_{\lambda,c}(z,w)=\left( \frac{2}{\lambda} \right)^{ \frac{1}{2\lambda} } (zw)^{\lambda-1} \widetilde{\kappa}_{\lambda,c} \left( \sqrt{ \frac{2}{\lambda} } z^\lambda, \sqrt{ \frac{2}{\lambda} } w^\lambda \right).
	\end{gather*}
	Here $\widetilde{\kappa}_{\lambda,c}$ is given in terms of a function $\widetilde{G}$ as
	\begin{gather*}
	\widetilde{\kappa}_{\lambda,c}(z,w)=\widetilde{G}(z,w)-\widetilde{G}(w,z)
	\end{gather*}
 and the following fractional differential equations hold:
	\begin{gather} \label{G zw FDE}
		D_z^{ \frac{1}{\lambda} } \widetilde{G}(z,w)= z^{ \frac{1}{\lambda} } \widetilde{G}(z,w)+(zw)^{\frac{1+c}{\lambda}-1} E_{\frac{2}{\lambda} \frac{1+c}{\lambda} } \big( (zw)^{\frac{2}{\lambda}} \big),
\\
			D_z^{ \frac{1}{\lambda} } \widetilde{\kappa}_{\lambda,c}(z,w)
			= z^{ \frac{1}{\lambda} } \widetilde{\kappa}_{\lambda,c}(z,w)+ (zw)^{ \frac{1+c}{\lambda}-1 } E_{ \frac{1}{\lambda},\frac{1+c}{\lambda} } \big( (zw)^{\frac{1}{\lambda}} \big)
			\nonumber\\
\hphantom{D_z^{ \frac{1}{\lambda} } \widetilde{\kappa}_{\lambda,c}(z,w)=}{}
 -\frac{ \Gamma\big(\frac{1+c}{\lambda}\big) }{\Gamma\big( \frac{c}{\lambda}\big) \Gamma\big( \frac{2+c}{\lambda}\big) } (zw)^{ \frac{1+c}{\lambda}-1 }(w/z)^{\frac{1}{\lambda}} E_{\frac{1}{\lambda},2,3+c-\lambda}\big(w^{\frac{2}{\lambda}}\big) .\label{FDE_bulkS}
		\end{gather}
\end{Proposition}


We call the equation \eqref{FDE_bulkS} (generalised) Christoffel--Darboux formula.
Such a differential equation satisfied by the correlation kernel, that makes the asymptotic analysis possible, is broadly called Christoffel--Darboux type formula, see, e.g., \cite{MR1917675,byun2021lemniscate,MR3450566}.

We also remark that the inhomogeneous term $(zw)^{ \frac{1+c}{\lambda}-1 } E_{ \frac{1}{\lambda},\frac{1+c}{\lambda} } \big( (zw)^{\frac{1}{\lambda}} \big)$ corresponds to the kernel of the complex Mittag-Leffler ensemble, see~\cite{ameur2018random}.
Such a relation has been observed for other models as well, which include the elliptic Ginibre ensemble~\cite{BE} and its chiral counter part~\cite{MR2180006}.
See also~\cite{MR1762659,MR1675356} for a similar relation for Hermitian ensembles, which gives an expression of the kernel of the symplectic ensemble in terms of a small number of orthogonal polynomials.

\begin{proof}[Proof of Proposition~\ref{Thm_bulkS}]
	Let us define
	\begin{gather*} %\label{kappa tilde0 bS}
		\widetilde{G}_{N}(z,w):=\sum_{k=0}^{N-1}z^{ \frac{2k+2+c}{\lambda}-1 } \prod_{j=0}^k \frac{ \Gamma\big( \frac{2j+1+c}{\lambda} \big) }{\Gamma\big( \frac{2j+2+c}{\lambda} \big)}
		\sum_{l=0}^{k} \frac{ w^{ \frac{2l+1+c}{\lambda}-1 } }{\Gamma\big( \frac{2l+1+c}{\lambda}\big)} \prod_{j=1}^l \frac{\Gamma\big( \frac{2j+c}{\lambda} \big)}{ \Gamma\big(\frac{2j-1+c}{\lambda}\big) } .
	\end{gather*}
	Here and henceforth we use the convention that if $l=0$, then $\prod_{j=1}^l \varphi=1$.
	By \eqref{SOP ML}, \eqref{sk ML} and~\eqref{micro scale ML}, the kernel $K_N$ is given by
	\begin{gather*}
			K_N(z,w)
			={\rm e}^{ -\frac{ |z|^{2\lambda}+|w|^{2\lambda} }{\lambda} } \begin{pmatrix}
				\kappa_N(z,w) & \kappa_N(z,\bar{w}) \\
				\kappa_N(\bar{z},w) & \kappa_N(\bar{z},\bar{w})
			\end{pmatrix},
		\end{gather*}
	where
	\begin{gather*}
		\kappa_N(z,w):=G_N(z,w)-G_N(w,z), \\ G_N(z,w):=\left( \frac{2}{\lambda} \right)^{\frac{1}{2\lambda}} (zw)^{\lambda-1} \widetilde{G}_N \left( \sqrt{ \frac{2}{\lambda} } z^\lambda, \sqrt{ \frac{2}{\lambda} } w^\lambda \right).
	\end{gather*}
	
	By definition of the Caputo derivative, we have that if $k>
	\lceil \nu \rceil$
	\begin{gather} \label{Caputo monomial}
		D_z^\nu z^k =\frac{\Gamma(k+1)}{\Gamma(k-\nu+1)}z^{k-\nu}
	\end{gather}
	and vanishes otherwise.
	Let us write $\widetilde{G}:=\lim\limits_{N \to \infty} \widetilde{G}_N,$ where the convergence is uniform on compact subsets of $\C$. The existence of the limit follows from the Weierstrass $M$-test.
	By direct computations using \eqref{Caputo monomial}, we have
	\begin{gather}
\widetilde{G}(z,w)-\frac{1}{\Gamma\big(\frac{2+c}{\lambda}\big)} z^{ \frac{2+c}{\lambda}-1 } w^{ \frac{1+c}{\lambda}-1 }
			\nonumber\\
\qquad{} =\sum_{k=1}^{\infty} z^{ \frac{2k+2+c}{\lambda}-1 } \prod_{j=0}^k \frac{ \Gamma\big( \frac{2j+1+c}{\lambda} \big) }{\Gamma\big( \frac{2j+2+c}{\lambda} \big)}
			\sum_{l=0}^{k} \frac{ w^{ \frac{2l+1+c}{\lambda}-1 } }{\Gamma\big( \frac{2l+1+c}{\lambda}\big)} \prod_{j=1}^l \frac{\Gamma\big( \frac{2j+c}{\lambda} \big)}{ \Gamma\big(\frac{2j-1+c}{\lambda}\big) } .\label{G zw til lam c 0}
		\end{gather}
	Note in particular that as $z\to 0,$
	\begin{gather} \label{G zw til asym}
		\widetilde{G}(z,w) \sim \frac{w^{ \frac{1+c}{\lambda}-1 } }{\Gamma\big(\frac{2+c}{\lambda}\big)} z^{ \frac{2+c}{\lambda}-1 }.
	\end{gather}
	Applying the operator $D_z^{ \frac{1}{\lambda} }$ to the identity \eqref{G zw til lam c 0}, we obtain
	\begin{gather*}
 D_z^{ \frac{1}{\lambda} } \widetilde{G}(z,w)-\frac{ 1}{\Gamma\big( \frac{1+c}{\lambda} \big)}(zw)^{ \frac{1+c}{\lambda}-1 }
		\\
\qquad{} = \sum_{k=1}^{\infty} z^{ \frac{2k+1+c}{\lambda}-1 } \prod_{j=1}^{k}\frac{ \Gamma \big( \frac{2j-1+c}{\lambda} \big) }{\Gamma\big( \frac{2j+c}{\lambda} \big)} \sum_{l=0}^{k} \frac{ w^{ \frac{2l+1+c}{\lambda}-1 } }{\Gamma\big( \frac{2l+1+c}{\lambda} \big)} \prod_{j=1}^l \frac{\Gamma\big(\frac{2j+c}{\lambda}\big) }{ \Gamma\big( \frac{2j-1+c}{\lambda} \big) }
		\\
\qquad{} = z^{ \frac{1}{\lambda} } \sum_{k=0}^{\infty} z^{ \frac{2k+2+c}{\lambda}-1 } \prod_{j=0}^k \frac{ \Gamma \big( \frac{2j+1+c}{\lambda} \big) }{\Gamma\big( \frac{2j+2+c}{\lambda} \big)} \sum_{l=0}^{k+1} \frac{ w^{ \frac{2l+1+c}{\lambda}-1 } }{\Gamma\big( \frac{2l+1+c}{\lambda} \big)} \prod_{j=1}^l \frac{\Gamma\big(\frac{2j+c}{\lambda}\big) }{ \Gamma\big( \frac{2j-1+c}{\lambda} \big) } ,
	\end{gather*}
	which leads to
	\begin{gather*}
		D_z^{ \frac{1}{\lambda} } \widetilde{G}(z,w)= z^{ \frac{1}{\lambda} } \widetilde{G}(z,w)
		+\sum_{k=0}^{\infty} \frac{ (zw)^{ \frac{2k+1+c}{\lambda}-1 } }{\Gamma\big( \frac{2k+1+c}{\lambda} \big)}.
	\end{gather*}
	Now \eqref{G zw FDE} follows from \eqref{ML function}.
	
	Similarly, by rearranging the terms, we have
	\begin{gather}
 \widetilde{G}(w,z)-\frac{ 1 }{\Gamma\big( \frac{1+c}{\lambda}\big)} z^{ \frac{1+c}{\lambda}-1 } \sum_{k=0}^{\infty} w^{ \frac{2k+2+c}{\lambda}-1 } \prod_{j=0}^k \frac{ \Gamma\big( \frac{2j+1+c}{\lambda} \big) }{\Gamma\big( \frac{2j+2+c}{\lambda} \big)}
			\nonumber\\
\qquad{} =\sum_{k=1}^{\infty} w^{ \frac{2k+2+c}{\lambda}-1 } \prod_{j=0}^k \frac{ \Gamma\big( \frac{2j+1+c}{\lambda} \big) }{\Gamma\big( \frac{2j+2+c}{\lambda} \big)} \sum_{l=1}^{k} \frac{ z^{ \frac{2l+1+c}{\lambda}-1 } }{\Gamma\big( \frac{2l+1+c}{\lambda} \big)} \prod_{j=1}^l \frac{\Gamma\big(\frac{2j+c}{\lambda}\big) }{ \Gamma\big( \frac{2j-1+c}{\lambda} \big) } .\label{G wz til lam c 0}
		\end{gather}
	Then by \eqref{ML3}, we have that as $z\to 0$,
	\begin{gather*} %\label{G wz til asym}
		\widetilde{G}(w,z)\sim \frac{1 }{ \Gamma\big( \frac{2+c}{\lambda}\big) } w^{ \frac{2+c}{\lambda}-1 } E_{\frac{1}{\lambda},2,3+c-\lambda}\big(w^{\frac{2}{\lambda}}\big) z^{ \frac{1+c}{\lambda}-1 }
	\end{gather*}
	Again, by applying the operator $D_z^{ \frac{1}{\lambda} }$ to \eqref{G wz til lam c 0}, we obtain
	\begin{gather*}
D_z^{ \frac{1}{\lambda} } \widetilde{G}(w,z)-\frac{ \Gamma\big(\frac{1+c}{\lambda}\big) }{\Gamma\big( \frac{c}{\lambda}\big) \Gamma\big( \frac{2+c}{\lambda}\big) } (zw)^{ \frac{1+c}{\lambda}-1 }(w/z)^{\frac{1}{\lambda}} E_{\frac{1}{\lambda},2,3+c-\lambda}\big(w^{\frac{2}{\lambda}}\big)\\
\qquad{} = z^{ \frac{1}{\lambda} } \sum_{k=1}^{\infty} w^{ \frac{2k+2+c}{\lambda}-1 } \prod_{j=0}^k \frac{ \Gamma\big( \frac{2j+1+c}{\lambda} \big) }{\Gamma\big( \frac{2j+2+c}{\lambda} \big)} \sum_{l=0}^{k-1} \frac{ z^{ \frac{2l+1+c}{\lambda}-1 } }{\Gamma\big( \frac{2l+1+c}{\lambda} \big)} \prod_{j=1}^l \frac{\Gamma\big(\frac{2j+c}{\lambda}\big) }{ \Gamma\big( \frac{2j-1+c}{\lambda} \big) }
		\\
\qquad{} = z^{ \frac{1}{\lambda} } \widetilde{G}(w,z)- \sum_{k=0}^{\infty} \frac{(zw)^{ \frac{2k+2+c}{\lambda}-1 }}{\Gamma\big( \frac{2k+2+c}{\lambda} \big)}.
	\end{gather*}
	Then by \eqref{ML function}, we have
	\begin{gather}
D_z^{ \frac{1}{\lambda} } \widetilde{G}(w,z)= z^{ \frac{1}{\lambda} } \widetilde{G}(w,z) -(zw)^{\frac{2+c}{\lambda}-1} E_{\frac{2}{\lambda},\frac{2+c}{\lambda}} \big((zw)^{\frac{2}{\lambda}}\big) \nonumber\\
\hphantom{D_z^{ \frac{1}{\lambda} } \widetilde{G}(w,z)=}{}
+\frac{ \Gamma(\frac{1+c}{\lambda}) }{\Gamma( \frac{c}{\lambda}) \Gamma( \frac{2+c}{\lambda}) } (zw)^{ \frac{1+c}{\lambda}-1 }(w/z)^{\frac{1}{\lambda}} E_{\frac{1}{\lambda},2,3+c-\lambda}\big(w^{\frac{2}{\lambda}}\big) . \label{G wz FDE}
		\end{gather}
	Now \eqref{FDE_bulkS} follows from \eqref{G zw FDE}, \eqref{G wz FDE} and \eqref{ML function}.
\end{proof}


\subsection[Bulk singularities of the order lambda=1/m with m in N]{Bulk singularities of the order $\boldsymbol{\lambda=\frac1m}$ with $\boldsymbol{m\in \mathbb{N}}$}

A point $p$ in which the eigenvalue density $\frac12\Delta Q(p)$ vanishes or diverges is called bulk singularity, see \cite{ameur2018random} and references therein. In this subsection, we shall consider the case $\lambda=\frac1m$, where $m$ is a positive integer and prove Theorem~\ref{Thm_ML integer} for general $c>-1$.

\begin{proof}[Proof of Theorem~\ref{Thm_ML integer}]
	By \eqref{G zw FDE} and \eqref{G zw til asym}, the function $\widetilde{G}$ is a unique solution to the $m$-th order linear (ordinary) differential equation
	\begin{gather} \label{G zw ODE}
		\partial_z ^m \widetilde{G}(z,w)=z^m \widetilde{G}(z,w)+F_{m,c}(zw),
	\end{gather}
	which satisfies the asymptotic behaviour
	\begin{gather} \label{G zw asym int}
		\widetilde{G}(z,w) \sim \frac{w^{m+mc-1}}{ \Gamma(2m+mc) } z^{2m+mc-1},\qquad z \to 0.
	\end{gather}
	Here the function $F_{m,c}$ is given by \eqref{F integer}. Note that the $m$ initial conditions are provided by differentiating \eqref{G zw asym int}.
	
We aim to solve the initial value problem \eqref{G zw ODE} and~\eqref{G zw asym int}.
	It is well known that the functions~$g_{j,m}$ given by \eqref{m indep sol} provide the solutions to the homogeneous equation of~\eqref{G zw ODE}, see, e.g., \cite[Theorem~5.18]{gorenflo2014mittag}.
	By~\eqref{ML3}, we have that
	\begin{gather} \label{gj asymp}
		g_{j,m}(z) \sim z^{j-1} ,\qquad z \to 0.
	\end{gather}
	Using this, it is easy to observe that
	\begin{gather*}
		W(g_{1,m},\dots,g_{m,m})(z)|_{z=0}=\prod_{j=0}^{m-1} j!=G(m+1),
	\end{gather*}
	where $G$ is the Barnes $G$-function \cite[Section~5.17]{olver2010nist}. Due to Abel's identity, for $m>1$, this leads to
	\begin{gather} \label{Wronskian g1m}
		W(g_{1,m},\dots,g_{m,m})(z) \equiv G(m+1).
	\end{gather}
	On the other hand for $m=1$, we have obviously $W(g_1)(z)=g_1(z).$
	
	For each $j \in \{1,\dots, m\}$, let us write
	\begin{gather*}
		\widetilde{W}_{j,m}(z,w):=\det
		\begin{bmatrix}
			g_{1,m}(z) & \dots & g_{j-1,m}(z) & 0 & g_{j+1,m}(z) & \dots & g_{m,m}(z)
			\\
			g_{1,m}'(z) & \dots & g_{j-1,m}'(z) & 0 & g_{j+1,m}'(z) & \dots & g_{m,m}'(z)
			\\
			\vdots & \vdots & \vdots & \vdots & \vdots & \vdots & \vdots
			\\
			g_{1,m}^{(m-1)}(z) & \dots & g_{j-1,m}^{(m-1)}(z) & F_{m,c}(zw) & g_{j+1,m}^{(m-1)}(z) & \dots & g_{m,m}^{(m-1)}(z)
		\end{bmatrix}.
	\end{gather*}
	Note here that by \eqref{W_j(z)}, we have
	\begin{gather} \label{Wj zw F Wj}
		\widetilde{W}_{j,m}(z,w)=(-1)^{m-j} F_{m,c}(zw) W_{j,m}(z).
	\end{gather}
	
	From now on, we shall only consider the case $m>1$. The other case $m=1$ follows along the same lines. In this case, one can also easily solve the associated linear non-homogenous ODE of degree $1$ employing an integrating factor.
	
	By virtue of the method of variation of parameters and \eqref{Wronskian g1m}, one can write
	\begin{align*}
			\widetilde{G}(z,w)&= \int_0^z \frac{1}{W(g_{1,m},\dots,g_{m,m})(s)}\sum_{j=1}^{m} 	\widetilde{W}_{j,m}(s,w) g_{j,m}(z)\,{\rm d}s+ \sum_{j=1}^{m} C_{j,m}(w) g_{j,m}(z)
			\\
			&=\frac{1}{G(m+1)}\int_0^z \sum_{j=1}^{m} 	\widetilde{W}_{j,m}(s,w) g_{j,m}(z)\,{\rm d}s+ \sum_{j=1}^{m} C_{j,m}(w) g_{j,m}(z)
	\end{align*}
	for some functions $C_{j,m}$. We shall show that in the above expression $\widetilde{G}$ satisfies the asymptotic behaviour~\eqref{G zw asym int} if and only if $C_{j,m} \equiv 0$ for every~$j$.
	
	By \eqref{gj asymp} and \eqref{Wronskian g1m}, direct computations of the determinant give that as $z \to 0$,
	\begin{gather*}
	W_{j,m}(z) \sim \frac{G(m+1)}{(m-1)!} \binom{m-1}{j-1} z^{m-j}.
	\end{gather*}
	On the other hand, by \eqref{ML function}, as $z \to 0$,
	\begin{gather*}
	F_{m,c}(zw) \sim \frac{w^{m+mc-1}}{\Gamma(m+mc)} z^{m+mc-1}.
	\end{gather*}
	Combining these two equations, we have that as $z\to 0$,
	\begin{gather*}
	\widetilde{W}_{j,m}(z,w) \sim \frac{w^{m+mc-1}}{\Gamma(m+mc)} (-1)^{m-j} \frac{G(m+1)}{(m-1)!} \binom{m-1}{j-1} z^{2m+mc-1-j}.
	\end{gather*}
	This leads to
	\begin{gather*}
	\frac{1}{G(m+1)}\int_0^z \widetilde{W}_{j,m}(s,w) g_{j,m}(z)\,{\rm d}s\\
\qquad{} \sim \frac{w^{m+mc-1}}{\Gamma(m+mc)} \frac{1}{(m-1)!} \binom{m-1}{j-1} \frac{(-1)^{m-j} }{2m+mc-j} z^{2m+mc-1}.
	\end{gather*}
	
	Next, we show the identity
	\begin{gather}\label{Comb identity}
		\frac{1}{(m-1)!} \sum_{j=1}^{m} \binom{m-1}{j-1} \frac{(-1)^{m-j} }{2m+mc-j} = \frac{\Gamma(m+mc)}{\Gamma(2m+mc)}.
	\end{gather}
	By the binomial theorem, we have
	\begin{gather*}
	\sum_{k=0}^{m-1} \binom{m-1}{k} x^{m+mc-1+k} = x^{m+mc-1}(1+x)^{m-1}.
	\end{gather*}
	Integrating this identity, we have
	\begin{gather*}
	\sum_{k=0}^{m-1} \binom{m-1}{k} \frac{x^{m+mc+k}}{m+mc+k} = \int_0^x t^{m+mc-1}(1+t)^{m-1}\,{\rm d}t.
	\end{gather*}
	Letting $x=-1$ and using the change of the variable $t=-s$, we have
	\begin{align*}
		\sum_{k=0}^{m-1} \binom{m-1}{k} \frac{(-1)^{k}}{m+mc+k} &= \int_0^1 s^{m+mc-1}(1-s)^{m-1}\,{\rm d}s
		\\
		&=B(m+mc,m)=(m-1)! \frac{\Gamma(m+mc)}{\Gamma(2m+mc)},
	\end{align*}
	where $B$ is the beta function. This gives~\eqref{Comb identity}.
	
	Using \eqref{Comb identity}, we have
	\begin{gather*}
	\frac{1}{G(m+1)}\int_0^z \sum_{j=1}^{m} 	\widetilde{W}_{j,m}(s,w) g_{j,m}(z)\,{\rm d}s \sim \frac{w^{m+mc-1}}{ \Gamma(2m+mc) } z^{2m+mc-1}.
	\end{gather*}
	Thus by \eqref{G zw asym int}, we obtain that $C_{j,m} \equiv 0$ for all $j \in \{1,\dots, m\}$.
	Therefore by~\eqref{Wj zw F Wj}, we conclude that
	\begin{align*}
		\widetilde{G}(z,w)&= \frac{1}{G(m+1)}\int_0^1 \sum_{j=1}^{m} 	\widetilde{W}_{j,m}(sz,w) z g_{j,m}(z)\,{\rm d}s
		\\
		&= \frac{1}{G(m+1)}\int_0^1 \sum_{j=1}^{m} (-1)^{m-j} 	\widetilde{W}_{j,m}(sz) z g_{j,m}(z)F_{m,c}(szw)\,{\rm d}s.
	\end{align*}
	This completes the proof.
\end{proof}
